%\usepackage{amsthm}
\usepackage{amsfonts}
\usepackage{amssymb}
%\usepackage{bbm}

\usepackage{dsfont}
\newcommand{\mathbbm}[1]{\mathds{#1}}  % keeps existing code working

\usepackage[normalem]{ulem}
%\usepackage{thmtools}
\usepackage{mathtools}
%\usepackage{thm-restate}
\usepackage{graphicx}
\usepackage{nicefrac}
%\usepackage{fullpage} <---- causes problems with margins
\usepackage{bm}
\usepackage{xcolor, xpatch}
\usepackage[linesnumbered,ruled]{algorithm2e}
\usepackage{tikz}
\usetikzlibrary{positioning, arrows.meta}
\usepackage{enumitem}
% provides \includegraphics{...} to
                              % include graphics (pdf format)
\usepackage{booktabs}         % fixes poor spacing in tables and
                              % provides \toprule, \midrule, \bottomrule
\usepackage[backend=bibtex,bibstyle=ieee,citestyle=numeric-comp,url=false,doi=false,isbn=false]{biblatex}
%\usepackage{algorithmicx}    % provides an algorithmic environment for
                              % describing algorithms. See
                              % https://ctan.org/pkg/algorithmicx
\usepackage[font=footnotesize,skip=10pt]{caption}
% \usepackage[caption=false,font=footnotesize]{subfig}
                              % provides subnumbering within a
                              % floating figure or table

%% For arrays and multiple-line equations, use the
%% IEEEeqnarray-environment. See
%%              https://moser-isi.ethz.ch/manuals.html#eqlatex  
%% for instructions.

%% Do NOT use amsthm or hyperref!
%% -IEEEtran provides its own versions of theorems.
%% -IEEEXplore does not accept submissions with hyperlinks


%%%%%
%% correct bad hyphenation here
\hyphenation{op-tical net-works semi-conduc-tor}

% \usepackage[draft]{hyperref}
% \hypersetup{
% 	colorlinks   = true, %Colours links instead of ugly boxes
% 	urlcolor     = blue, %Colour for external hyperlinks
% 	linkcolor    = blue, %Colour of internal links
% 	citecolor   = red %Colour of citations
% }
\usepackage{pgfplots}
\DeclareMathOperator*{\argmax}{arg\,max}
\DeclareMathOperator*{\argmin}{arg\,min}


\makeatletter
\xpatchcmd{\proof}{\topsep6\p@\@plus6\p@\relax}{}{}{}
\makeatother


% \newtheorem{thm}{theorem}[section]
% \newtheorem{prop}[thm]{Proposition}
% \newtheorem{lemma}[thm]{Lemma}
% \newtheorem{theorem}[thm]{Theorem}
% \newtheorem{remark}[thm]{Remark}
% \newtheorem{cor}[thm]{Corollary}
% \newtheorem{con}[thm]{Condition}
% \newtheorem{claim}[thm]{Claim}
% \newtheorem{definition}[thm]{Definition}
% \newtheorem{cnst}[thm]{Construction}

% \crefname{section}{Section}{Sections}
% \crefname{subsection}{Section}{Sections}
% \crefname{equation}{Eq.}{Equations}
% \crefname{enumi}{part}{parts}
% \crefname{table}{Table}{Tables}
% \crefname{figure}{Figure}{Figures}
% \crefname{algocf}{Algorithm}{Algorithms}

\usepackage{amsthm}  
\usepackage{aliascnt} % allows for a single counter in \autoref'ed theorems
\usepackage[draft]{hyperref} % cross-reference within document. Options: [draft,hidelinks,linktocpage]
%\usepackage[hidelinks,linktocpage]{hyperref}
\usepackage{cleveref} % enhances LaTeX's cross-referencing features
%\usepackage{upref}
%\usepackage{mathrsfs}

\newtheorem{theorem}{Theorem}
\crefname{theorem}{Theorem}{Theorems}

% \theoremheaderfont{\normalfont\bfseries}
% \theorembodyfont{\normalfont}
\theoremstyle{plain}

\newaliascnt{lemma}{theorem}
\newtheorem{lemma}[lemma]{Lemma}
\aliascntresetthe{lemma}
\crefname{lemma}{Lemma}{Lemmas}

\newaliascnt{definition}{theorem}
\newtheorem{definition}[definition]{Definition}
\aliascntresetthe{definition}
\crefname{definition}{Definition}{Definitions}

\newaliascnt{corollary}{theorem}
\newtheorem{corollary}[corollary]{Corollary}
\aliascntresetthe{corollary}
\crefname{corollary}{Corollary}{Corollaries}

\newaliascnt{claim}{theorem}
\newtheorem{claim}[claim]{Claim}
\aliascntresetthe{claim}
\crefname{claim}{Claim}{Claims}

\newaliascnt{conjecture}{theorem}
\newtheorem{conjecture}[conjecture]{Conjecture}
\aliascntresetthe{conjecture}
\crefname{conjecture}{Conjecture}{Conjectures}

\newaliascnt{question}{theorem}
\newtheorem{question}[question]{Question}
\aliascntresetthe{question}
\crefname{question}{Question}{Questions}

\newaliascnt{problem}{theorem}
\newtheorem{problem}[problem]{Problem}
\aliascntresetthe{problem}
\crefname{problem}{Problem}{Problems}

\newaliascnt{example}{theorem}
\newtheorem{example}[example]{Example}
\aliascntresetthe{example}
\crefname{example}{Example}{Examples}

\newaliascnt{oquestion}{theorem}
\newtheorem{oquestion}[oquestion]{Open Question}
\aliascntresetthe{oquestion}
\crefname{oquestion}{Open Question}{Open Questions}

\newaliascnt{remark}{theorem}
\newtheorem{remark}[remark]{Remark}
\aliascntresetthe{remark}
\crefname{remark}{Remark}{Remarks}



%\theoremstyle{remark}


% \newtheorem*{lemma*}{Lemma}
% \newtheorem*{theorem*}{theorem}

% ------- Math environments

% \DeclarePairedDelimiter{\ceil}{\lceil}{\rceil}
% \DeclarePairedDelimiter{\floor}{\lfloor}{\rfloor}

% \DeclarePairedDelimiter\abs{\lvert}{\rvert}
% %\DeclarePairedDelimiter\ceilenv{\lceil}{\rceil}
% %\DeclarePairedDelimiter\floorenv{\lfloor}{\rfloor}
% \DeclarePairedDelimiter\parenv{\lparen}{\rparen}
% \DeclarePairedDelimiter\sparenv{\lbrack}{\rbrack}
% \DeclarePairedDelimiter\bracenv{\lbrace}{\rbrace}
% \DeclarePairedDelimiterX\mathset[2]{\lbrace}{\rbrace}{#1 : #2}


% ------- Calligraphy

\newcommand{\bfa}{{\boldsymbol a}}
\newcommand{\bfb}{{\boldsymbol b}}
\newcommand{\bfc}{{\boldsymbol c}}
\newcommand{\bfd}{{\boldsymbol d}}
\newcommand{\bfe}{{\boldsymbol e}}
\newcommand{\bff}{{\boldsymbol f}}
\newcommand{\bfg}{{\boldsymbol g}}
\newcommand{\bfh}{{\boldsymbol h}}
\newcommand{\bfi}{{\boldsymbol i}}
\newcommand{\bfj}{{\boldsymbol j}}
\newcommand{\bfk}{{\boldsymbol k}}
\newcommand{\bfl}{{\boldsymbol l}}
\newcommand{\bfm}{{\boldsymbol m}}
\newcommand{\bfn}{{\boldsymbol n}}
\newcommand{\bfo}{{\boldsymbol o}}
\newcommand{\bfp}{{\boldsymbol p}}
\newcommand{\bfq}{{\boldsymbol q}}
\newcommand{\bfr}{{\boldsymbol r}}
\newcommand{\bfs}{{\boldsymbol s}}
\newcommand{\bft}{{\boldsymbol t}}
\newcommand{\bfu}{{\boldsymbol u}}
\newcommand{\bfv}{{\boldsymbol v}}
\newcommand{\bfw}{{\boldsymbol w}}
\newcommand{\bfx}{{\boldsymbol x}}
\newcommand{\bfy}{{\boldsymbol y}}
\newcommand{\bfz}{{\boldsymbol z}}
\newcommand{\bfA}{{\mathbf A}}
\newcommand{\bfB}{{\mathbf B}}
\newcommand{\bfC}{{\mathbf C}}
\newcommand{\bfD}{{\mathbf D}}
\newcommand{\bfE}{{\mathbf E}}
\newcommand{\bfF}{{\mathbf F}}
\newcommand{\bfG}{{\mathbf G}}
\newcommand{\bfH}{{\mathbf H}}
\newcommand{\bfI}{{\mathbf I}}
\newcommand{\bfJ}{{\mathbf J}}
\newcommand{\bfK}{{\mathbf K}}
\newcommand{\bfL}{{\mathbf L}}
\newcommand{\bfM}{{\mathbf M}}
\newcommand{\bfN}{{\mathbf N}}
\newcommand{\bfO}{{\mathbf O}}
\newcommand{\bfP}{{\mathbf P}}
\newcommand{\bfQ}{{\mathbf Q}}
\newcommand{\bfR}{{\mathbf R}}
\newcommand{\bfS}{{\mathbf S}}
\newcommand{\bfT}{{\mathbf T}}
\newcommand{\bfU}{{\mathbf U}}
\newcommand{\bfV}{{\mathbf V}}
\newcommand{\bfW}{{\mathbf W}}
\newcommand{\bfX}{{\mathbf X}}
\newcommand{\bfY}{{\mathbf Y}}
\newcommand{\bfZ}{{\mathbf Z}}
\newcommand{\bfzero}{{\mathbf 0}}
\newcommand{\bfone}{{\mathbf 1}}

\newcommand{\cA}{\mathcal{A}}
\newcommand{\cB}{\mathcal{B}}
\newcommand{\cC}{\mathcal{C}}
\newcommand{\cD}{\mathcal{D}}
\newcommand{\cE}{\mathcal{E}}
\newcommand{\cF}{\mathcal{F}}
\newcommand{\cG}{\mathcal{G}}
\newcommand{\cH}{\mathcal{H}}
\newcommand{\cI}{\mathcal{I}}
\newcommand{\cJ}{\mathcal{J}}
\newcommand{\cK}{\mathcal{K}}
\newcommand{\cL}{\mathcal{L}}
\newcommand{\cM}{\mathcal{M}}
\newcommand{\cN}{\mathcal{N}}
\newcommand{\cO}{\mathcal{O}}
\newcommand{\cP}{\mathcal{P}}
\newcommand{\cQ}{\mathcal{Q}}
\newcommand{\cR}{\mathcal{R}}
\newcommand{\cS}{\mathcal{S}}
\newcommand{\cT}{\mathcal{T}}
\newcommand{\cU}{\mathcal{U}}
\newcommand{\cV}{\mathcal{V}}
\newcommand{\cW}{\mathcal{W}}
\newcommand{\cX}{\mathcal{X}}
\newcommand{\cY}{\mathcal{Y}}
\newcommand{\cZ}{\mathcal{Z}}

% ------- Math environments

\DeclarePairedDelimiter\abs{\lvert}{\rvert}
\DeclarePairedDelimiter\ceilenv{\lceil}{\rceil}
\DeclarePairedDelimiter\floorenv{\lfloor}{\rfloor}
\DeclarePairedDelimiter\parenv{\lparen}{\rparen}
\DeclarePairedDelimiter\sparenv{\lbrack}{\rbrack}
\DeclarePairedDelimiter\bracenv{\lbrace}{\rbrace}
\DeclarePairedDelimiter\multbrace{\lbrace\!\!\lbrace}{\rbrace\!\!\rbrace}
%\DeclarePairedDelimiterX\mathset[2]{\lbrace}{\rbrace}{#1 \mathrel{}\delimsize\vert\mathrel{} #2}
\DeclarePairedDelimiterX\mathset[2]{\lbrace}{\rbrace}{#1 : #2}
\DeclarePairedDelimiterX\multset[2]{\lbrace\!\!\lbrace}{\rbrace\!\!\rbrace}{#1 : #2}


% ------- New Math symbols
\newcommand{\deq}{\mathrel{\triangleq}}



% ------- New Math Operators
\newcommand{\floor}[1]{\lfloor #1 \rfloor}
\newcommand{\ceil}[1]{\left \lceil #1 \right \rceil}
\newcommand{\Fq}{\mathbb{F}_q}
\newcommand{\Fp}{\mathbb{F}_p}
\newcommand{\Fqm}[1]{\mathbb{F}_{q^{#1}}}
\newcommand{\Ftwo}{\mathbb{F}_2}
\newcommand{\dham}{d_{\mathrm{H}}}
\newcommand{\ded}{d_{\textup{ed}}}
\newcommand{\dlev}{d_{\mathrm{L}}}
\newcommand{\ballins}[2]{\cI_{#2}(#1)}
\newcommand{\balldel}[2]{\cD_{#2}(#1)}
\newcommand{\balldelp}[2]{\cD'_{#2}(#1)}
\newcommand{\ballinsp}[2]{\cI'_{#2}(#1)}

\newcommand{\balpha}{{\boldsymbol\alpha}}
\newcommand{\prob}{\mathrm{Pr}}
\newcommand{\ex}{\mathbb{E}}


\newcommand{\vt}[2]{\mathcal{VT}_{#1}(#2)}
\newcommand{\tauD}{\tau_{\mathrm{D}}}
\newcommand{\tauI}{\tau_{\mathrm{I}}}
\newcommand{\indicator}{\mathbbm{1}}
\newcommand{\listsub}[2]{\bar{L}_\mathrm{sub}(#1,#2)}
\newcommand{\listins}[2]{\bar{L}_\mathrm{ins}(#1,#2)}
\newcommand{\listdel}[2]{\bar{L}_\mathrm{del}(#1,#2)}
%\renewcommand{\listdel}[2]{\smash[b]{\listdel{#1}{#2}}}



% --------- Davey-MacKay channel params
\newcommand{\pin}{P_{\mathrm{i}}}
\newcommand{\pd}{P_{\mathrm{d}}}
\newcommand{\ps}{P_{\mathrm{s}}}
\newcommand{\pt}{P_{\mathrm{t}}}
\newcommand{\tr}{\mathtt{Tr}}
\newcommand{\del}{\mathtt{Del}}


% -------------- Comments

\newcommand{\rc}[1]{{\color{blue}[Roni: #1]}}
\newcommand{\ab}[1]{{\color{red}{[AB: #1]}}}
\newcommand{\aw}[1]{\textcolor{magenta}{[Antonia: #1]}}
\newcommand{\ey}[1]{{\color{red}[Eitan: #1]}}

\bibliography{resources/literature}
\renewcommand*{\bibfont}{\footnotesize}

\usepackage{fancyhdr}
\renewcommand{\headrulewidth}{0pt}